\documentclass[12pt]{article}
\usepackage{microtype}
\usepackage{a4wide}
\usepackage{amsmath, amssymb, amsthm}
\usepackage{tikz-cd}
\newtheorem{theorem}{Theorem}
\newtheorem{lemma}{Lemma}
\newtheorem{proposition}{Proposition}
\newtheorem{corollary}{Corollary}
\newtheorem{definition}{Definition}

\newcommand{\cat}{%
\mathbf%
}

\newcommand{\idarrow}[1][]{%
\mathbf{1}_{#1}%
}

\title{Category theory}
\begin{document}

\section{Category}

\newcommand{\morphs}[3]{\cat{#1}_{#2}^{#3}}

\begin{definition}
    A category \(\cat{C}\) consists of 
    \begin{itemize}
        \item A class \(ob~\cat{C}\), whose elements are called objects, we note \(a : \cat{C}\) if
            a is an element of \(ob~\cat{C}\).
        \item A family of classes \(hom~\cat{C}~{a}~{b}\) indexed by pairs of objects \(a\) and \(b\), 
            whose elements are called morphisms from \(a\) to \(b\), we note \(f : \morphs{C}{a}{b}\)
            if \(f\) is an element of \(hom~\cat{C}~{a}~{b}\).
        \item A family of identity morphisms \(\idarrow[a]\) indexed by objects \(a\)
        \item A binary operation \(\circ\), called composition of morphisms, such that 
            for any three objects \(a\), \(b\) and \(c\) we have 
            \begin{equation*}
                \circ : \morphs{C}{b}{c} \times \morphs{C}{a}{b} \rightarrow \morphs{C}{a}{c}    
            \end{equation*}
            such that 
            \begin{itemize}
                \item for every \(f : \morphs{C}{a}{b}\) we have
                \(\idarrow[b] \circ f = f \circ \idarrow[a] = f\)
                \begin{equation*}
                    \begin{tikzcd}
                        a   \arrow[r,"f"]
                            \arrow[loop,out=150,in=210,distance=20,swap,"{\mathbf{1}}_a"]
                            \arrow[r,bend left =70,"f\circ{\mathbf{1}}_a"] 
                            \arrow[r, bend right = 70,swap,"{\mathbf{1}}_b\circ f"]&
                        b   \arrow[loop,out=30,in=-30,distance=20,"{\mathbf{1}}_b"]
                    \end{tikzcd}
                \end{equation*}
                \item $\circ$ is associative, i.e. for every \(f : \morphs{C}{a}{b}\), 
                \(g : \morphs{C}{b}{c}\) and \(h : \morphs{C}{c}{d}\) we have
                \(h \circ (g \circ f) = (h \circ g) \circ f\)
                \begin{equation*}
                    \begin{tikzcd}
                        a   \arrow[r,"f"]
                            \arrow[dr, swap,"g \circ f"]& 
                        b   \arrow[dr,"g \circ h"] 
                            \arrow[d,swap,"g"]\\
                        {} & 
                        c \arrow[r,swap,"h"] & 
                        d
                    \end{tikzcd}
                \end{equation*}
                
            \end{itemize}
    \end{itemize}
\end{definition}

\section{Properties of morphisms}

A monomorphism is a left-cancellative morphism.
\begin{definition}
    Let \(\cat{C}\) be a category and \(a,b:\cat{C}\), a morphism \(f : \morphs{C} {b} {c}\) is a monomorphism
    if, for every \(g, h : \morphs{C}{a}{b}\), the equality \(f \circ g = f \circ h\) implies that \(g=h\). 
\end{definition}

\begin{definition}
    Let \(\cat{C}\) be a category and \(a,b:\cat{C}\), a morphism \(f : \morphs{C} {a} {b}\) is an epimorphism
    if, for every \(g, h : \morphs{C}{b}{c}\), the equality \(g \circ f = h \circ f \) implies that \(g=h\). 
\end{definition}

\begin{definition}
    Let \(\cat{C}\) be a category and \(a,b:\cat{C}\), a morphism \(f : \morphs{C} {a} {b}\) is an isomorphism if there exists \(f^{-1} : \morphs{C}{b}{a}\), called the inverse of \(f\), such that \(f^{-1} \circ f = \idarrow[a]\) and 
    \(f \circ f^{-1} = \idarrow[b]\).
\end{definition}

In general, a morphism that is both a monomorphism and an epimorphism is not a isomorphism.

A retraction is the left-inverse of a morphism.
\begin{definition}
    Let \(\cat{C}\) be a category, \(a,b:\cat{C}\) and \(f : \morphs{C}{a}{b}\),
    a morphism \( g : \morphs{C}{b}{a}\) is a retraction of \(f\) if 
    \(g \circ f = \idarrow[a]\).
\end{definition}
A section is the right-inverse of a morphism.
\begin{definition}
    Let \(\cat{C}\) be a category, \(a,b:\cat{C}\) and \(f : \morphs{C} {a} {b}\),
    a morphism \( g : \morphs{C}{b}{a}\) is a section of \(f\) if 
    \(f \circ g = \idarrow[b]\).
\end{definition}

\begin{lemma}
    Let \(\cat{C}\) be a category, \(a,b:\cat{C}\),
    if \(g : \morphs{C}{b}{a}\) is a section of some \(f :\morphs{C}{a}{b}\) then it is also a monomorphism.
\end{lemma}

\begin{lemma}
    Let \(\cat{C}\) be a category, \(a,b:\cat{C}\) and \(f : \morphs{C} {a} {b}\),
    if \(g : \morphs{C} {b} {a}\) is a retraction of some \(f :\morphs{C}{a}{b}\) then it is also an epimorphism.
\end{lemma}

\begin{lemma}
    Let \(\cat{C}\) be a category and \(b c:\cat{C}\), \(f : \morphs{C} {b} {c}\) 
    is an isomorphism then it is also an monomorphism.
\end{lemma}
\begin{proof}
    By definition of a monomorphism, we must prove that, for all \(a:\cat{C}\) and for all \(g, h:\morphs{C}{a}{b}\),
    the equality \(f \circ g = f \circ h\) implies \(g = h\).
    By the hypothesis on \(f\), there exists an inverse \(f^{-1}\) such that \(f^{-1}\circ f = \idarrow[a]\), Using the properties of composition, we can express \(g\) as \(f^{-1} \circ (f \circ g)\).
    From the equality hypothesis, we have \(g = f^{-1} \circ (f \circ h)\).
    Finally, using the properties of composition again, we conclude that $g = h$.
\end{proof}

\begin{lemma}
    Let \(\cat{C}\) be a category and \(a,b:\cat{C}\), \(f : \morphs{C} {a} {b}\) 
    is an isomorphism then it is also an epimorphism.
\end{lemma}
\begin{proof}
By definition of an epimorphism, we must prove that, for all \(c:\cat{C}\) and for all \(g, h:\morphs{C}{b}{c}\),
the equality \(g \circ f = h \circ f\) implies \(g = h\).
By the hypothesis on \(f\), there exists an inverse \(f^{-1}\) such that \(f\circ f^{-1} = \idarrow[b]\), Using the properties of composition, we can express \(g\) as \(g \circ f \circ f^{-1}\).
From the equality hypothesis, we have \(g = h \circ f \circ f^{-1}\).
Finally, using the properties of composition again, we conclude that $g = h$.
\end{proof}
\begin{lemma} 
    Let \(\cat{C}\) be a category and \(a,b:\cat{C}\), if \(g : \morphs{C} {b} {b
    a}\) is both a retraction and a monomorphism then it is also an isomorphism. 
\end{lemma}
\begin{proof}
    By definition of an isomorphism we must find a morphism \(g^{-1} : \morphs{C}{a}{b}\) such that \(g \circ g^{-1} = \idarrow[a]\) and \(g^{-1} \circ g = \idarrow[b]\).
    By hypothesis, \(g\) is a retraction, thus there exists \(f : \morphs{C}{a}{b}\) such that the equation \(g \circ f = \idarrow[a]\) holds. To prove that \(g\) is an isomorphism, it suffice to prove that \(f \circ g = \idarrow[b]\). Using the properties of composition and the equation we have   
    \( g \circ (f \circ g) = g \circ \idarrow[b]\).
    But, by hypothesis, \(g\) is also a monorphism and thus we have
    \((f \circ g) = \idarrow[b]\).
\end{proof}

\begin{lemma} 
    Let \(\cat{C}\) be a category and \(a,b:\cat{C}\), if \(g : \morphs{C} {b} {a}\) is both a section and an epimorphism then it is also an isomorphism. 
\end{lemma}
\begin{proof}
    By definition of an isomorphism we must find a morphism \(g^{-1} : \morphs{C}{a}{b}\) such that \(g \circ g^{-1} = \idarrow[b]\) and \(g^{-1} \circ g = \idarrow[a]\).
    By hypothesis, \(g\) is a section, thus there exists \(f : \morphs{C}{a}{b}\) such that the equation \(f \circ g = \idarrow[b]\) holds. To prove that \(g\) is an isomorphism, it suffice to prove that \(g \circ f = \idarrow[a]\). Using the properties of composition and the equation we have   
    \( g \circ f \circ g = \idarrow[a] \circ g\).
    But, by hypothesis, \(g\) is also a monorphism and thus we have
    \((g \circ f) = \idarrow[a]\).
\end{proof}


\section{Functors}

\begin{definition}
    Given two categories \(\cat{C}\) and \(\cat{D}\), a functor from \(\cat{C}\) to \(\cat{D}\) is a mapping that 
    \begin{itemize}
        \item associates each \(a : \cat{C}\) to an object \(F a : \cat{D}\)
        \item associates each \(f : \morphs{C}{a}{b}\) to a morphism \(F f : \morphs{D}{F a}{F b}\)
        such that the following conditions hold 
        \begin{itemize}
            \item \(F \idarrow[a] = \idarrow[F a]\)
            \item \(F (g \circ f) = F g \circ F f\)
        \end{itemize}
        \begin{equation*}
            \begin{tikzcd}
                F a \arrow[r,"F f"]
                    \arrow[rr, bend left=70,"F (g\circ f)"]
                    \arrow[rr, bend right=70,"(F g) \circ (F f)"] 
                &
                F b \arrow[r,"F g"]
                & 
                F c
            \end{tikzcd}
        \end{equation*}
    \end{itemize}
\end{definition}

\begin{paragraph}{Identity functor}
Given a category \(\cat{C}\), we define the identity mapping \(\idarrow[\cat{C}]\) by
\(\idarrow[\cat{C}] a = a\) and \(\idarrow[\cat{C}] f = f\) for all \(a:\cat{C}\)
and \(f:\morphs{C}{a}{b}\). Obviously, \(\idarrow[\cat{C}]\) is a functor
as by definition we have 
\(\idarrow[\cat{C}] \idarrow[a] = \idarrow[\cat{C} a]\) and 
\(\idarrow[\cat{C}] (g \circ f) = \idarrow[\cat{C}]g\circ\idarrow[\cat{C}] f\)
\end{paragraph}
\begin{paragraph}{Composition of functors}
We define a composition operator $\circ$ that, 
given three categories \(\cat{C}\), \(\cat{D}\) and \(\cat{E}\), associates
two functors $G$ from \(\cat{D}\) to \(\cat{E}\) and $F$ from \(\cat{C}\) and \(\cat{D}\) to the functor \(G \circ F\) defined by 
\((G \circ F) a = G (F a)\) and \((G \circ F) f = G (F f)\) for all \(a:\cat{C}\)
and \(f:\morphs{C}{a}{b}\).
Obviously \(G \circ F\) is a functor as from properties of $G$ and $F$ we have
\begin{itemize}
\item \((G \circ F) \idarrow[a] = G(\idarrow[F a]) = \idarrow[(G\circ F)a]\)
\item \((G \circ F)(g \circ f) = 
    G (F g \circ F f) = (G \circ F) g \circ (G \circ F) f\)
\end{itemize}
\end{paragraph}


The category \(\cat{Cat}\) is the category whose objects are small categories
and morphisms from \(C:\cat{Cat}\) to \(D:\cat{Cat}\) are functors from
from \(C:\cat{Cat}\) to \(D:\cat{Cat}\).
The identify functor \(\idarrow[C]\) from \(C\) to \(C\) and the composition
\((G \circ F)\) from \(C\) to \(E\) of two functors \(F\) from \(C\) to \(D\) and \(G\) from \(D\) to \(E\) are defined by
\begin{itemize}
    \item \(\idarrow[C] a = a\) for all \(a : \cat{C}\)
    \item \(\idarrow[C] f = f\) for all \(f : \morphs{C}{a}{b}\)
    \item \((G \circ F) a = G (F a)\) for all \(a : \cat{C}\) 
    \item \((G \circ F) f = G (F f)\) for all \(f : \morphs{C}{a}{b}\)
\end{itemize}

\begin{proposition}
    Small categories and functors together with the composition of functors 
    form a category
\end{proposition}
\begin{proof}
\begin{itemize}
    \item left and right identities 
    \item associativity
\end{itemize}
\end{proof}

\section{Category of small categories}

The category \(\mathbf{Cat}\) of small categories is the category whose objects are small categories
and morphisms are functors.

\begin{itemize}
    \item The identity functor is defined by \(\idarrow[C] x = x\) and \(\idarrow[C] f = f\)
    \item The composition operator is defined by \((G \circ F) a = G(F a)\) and \((G \circ F) f = G(F f)\)
\end{itemize}
We have to prove that 
\begin{itemize}
    \item 
\end{itemize}

initial object empty, the terminal is the singleton.
\end{document}